% Compile with: uplatex file.tex && dvipdfmx file.dvi
%           or: latexmk -uplatex file.tex
%\documentclass[uplatex,dvipdfmx]{jsarticle}
\documentclass[dvipdfmx,a4paper,11pt]{jsarticle}
\usepackage{amsmath, amssymb}
\usepackage{tikz}

\title{scenario\_coulomb\_gauge}
\date{\today}

\begin{document}
\maketitle

スカラーポテンシャル $\phi$（空間 $x,y,z$ と時間 $t$ に依存）と、ベクトルポテンシャル $\mathbf A$（成分 $A_x,A_y,A_z$、これも時間に依存）を宣言する。この2つから電場 $\mathbf E$・磁場 $\mathbf B$ を組み立てる。

phi を時間依存のスカラー場として定義する

\[
  \phi
\]

A をベクトル場として定義する

\[
  A
\]

$A_x$ を時間依存関数として定義する

$\checkmark$ A\_x を時間依存として宣言しました（時間変数: t）

$A_y$ を時間依存関数として定義する

$\checkmark$ A\_y を時間依存として宣言しました（時間変数: t）

$A_z$ を時間依存関数として定義する

$\checkmark$ A\_z を時間依存として宣言しました（時間変数: t）

後で偏微分の計算に使うため、$\phi,A_x,A_y,A_z$ それぞれの値に別名（phi0, Ax0, Ay0, Az0）をつけて式番号として登録しておく。

$\phi_{0} = \phi$

\begin{equation}
  \phi_{0} = \phi
  \label{eq:E1}
\end{equation}

$Ax0 = A_x$

\begin{equation}
  Ax0 = A_x
  \label{eq:E2}
\end{equation}

$Ay0 = A_y$

\begin{equation}
  Ay0 = A_y
  \label{eq:E3}
\end{equation}

$Az0 = A_z$

\begin{equation}
  Az0 = A_z
  \label{eq:E4}
\end{equation}

電場の定義 $\mathbf E=-\nabla\phi-\partial\mathbf A/\partial t$ の1項目 $\nabla\phi$ に使うため、$\phi$ の空間偏微分 $\partial\phi/\partial x,\partial\phi/\partial y,\partial\phi/\partial z$ を求める。

E1 を x で偏微分して

\begin{equation}
  \frac{\partial \phi_{0}}{\partial x} = \frac{d}{d x} \phi
  \label{eq:E5}
\end{equation}

E1 を y で偏微分して

\begin{equation}
  \frac{\partial \phi_{0}}{\partial y} = \frac{d}{d y} \phi
  \label{eq:E6}
\end{equation}

E1 を z で偏微分して

\begin{equation}
  \frac{\partial \phi_{0}}{\partial z} = \frac{d}{d z} \phi
  \label{eq:E7}
\end{equation}

電場の定義の2項目 $\partial\mathbf A/\partial t$ に使うため、$\mathbf A$ の各成分の時間微分 $\partial A_x/\partial t,\partial A_y/\partial t,\partial A_z/\partial t$ を求める。

E2 を t で偏微分して

\begin{equation}
  \frac{\partial Ax0}{\partial t} = \frac{d}{d t} A_{x}
  \label{eq:E8}
\end{equation}

E3 を t で偏微分して

\begin{equation}
  \frac{\partial Ay0}{\partial t} = \frac{d}{d t} A_{y}
  \label{eq:E9}
\end{equation}

E4 を t で偏微分して

\begin{equation}
  \frac{\partial Az0}{\partial t} = \frac{d}{d t} A_{z}
  \label{eq:E10}
\end{equation}

上で求めた偏微分を組み合わせ、電場 $\mathbf E=-\nabla\phi-\partial\mathbf A/\partial t$ の各成分を組み立てる。

$E_x = -E5 - E8$

\begin{equation}
  E_x = -\left(\frac{d}{d x} \phi\right)-\frac{d}{d t} A_{x}
  \label{eq:E11}
\end{equation}

$E_y = -E6 - E9$

\begin{equation}
  E_y = -\left(\frac{d}{d y} \phi\right)-\frac{d}{d t} A_{y}
  \label{eq:E12}
\end{equation}

$E_z = -E7 - E10$

\begin{equation}
  E_z = -\left(\frac{d}{d z} \phi\right)-\frac{d}{d t} A_{z}
  \label{eq:E13}
\end{equation}

磁場 $\mathbf B=\nabla\times\mathbf A$ を、$\nabla\times$ 演算子を $\mathbf A$ に直接適用して1ステップで求める（3成分がまとめて返る）。

$\nabla \times A$

\begin{equation}
  \left[ \frac{d}{d y} A_{z}-\frac{d}{d z} A_{y} , \frac{d}{d z} A_{x}-\frac{d}{d x} A_{z} , \frac{d}{d x} A_{y}-\frac{d}{d y} A_{x} \right]
  \label{eq:E14}
\end{equation}

B をベクトル場として定義する

\[
  B
\]

得られた curl の3成分を、ベクトル場 $\mathbf B$ の $x,y,z$ 成分としてそれぞれ束縛する。

$B_x = E14[1]$

\begin{equation}
  B_x = \frac{d}{d y} A_{z}-\frac{d}{d z} A_{y}
  \label{eq:E15}
\end{equation}

$B_y = E14[2]$

\begin{equation}
  B_y = \frac{d}{d z} A_{x}-\frac{d}{d x} A_{z}
  \label{eq:E16}
\end{equation}

$B_z = E14[3]$

\begin{equation}
  B_z = \frac{d}{d x} A_{y}-\frac{d}{d y} A_{x}
  \label{eq:E17}
\end{equation}

$\mathbf B=\nabla\times\mathbf A$ から常に成り立つはずの恒等式 $\nabla\cdot\mathbf B=0$（磁荷が存在しないこと）を、$\mathbf B$ の発散を計算して確認する。

$\nabla \cdot B$

\begin{equation}
  0
  \label{eq:E18}
\end{equation}

クーロンゲージ条件 $\nabla\cdot\mathbf A=0$ を確認するため、まず $\mathbf A$ の発散を計算する。

$\nabla \cdot A$

\begin{equation}
  \frac{d}{d z} A_{z}+\frac{d}{d y} A_{y}+\frac{d}{d x} A_{x}
  \label{eq:E19}
\end{equation}

この発散の値に gauge という名前をつける。

$gauge = E19$

\begin{equation}
  gauge = \frac{d}{d z} A_{z}+\frac{d}{d y} A_{y}+\frac{d}{d x} A_{x}
  \label{eq:E20}
\end{equation}

クーロンゲージでは $\nabla\cdot\mathbf A=0$ という条件を課す。gauge の値を 0 と置いた方程式として確定する。

E20 の右辺を 0 とする

\begin{equation}
  \frac{d}{d z} A_{z}+\frac{d}{d y} A_{y}+\frac{d}{d x} A_{x} = 0
  \label{eq:E21}
\end{equation}

後でこのゲージ条件を解くために、$A_z$ の空間微分 $\partial A_z/\partial z$ を個別に用意し、名前をつける。

Az0 を z で偏微分して

\begin{equation}
  \frac{\partial Az0}{\partial z} = \frac{d}{d z} A_{z}
  \label{eq:E22}
\end{equation}

$Az_z = E22$

\begin{equation}
  Az_z = \frac{d}{d z} A_{z}
  \label{eq:E23}
\end{equation}

ゲージ条件の式を $\partial A_z/\partial z$ について解き、$\partial A_x/\partial x,\partial A_y/\partial y$ で表す。

E21 を $Az_z$ について解いて

\begin{equation}
  \frac{d}{d z} A_{z} = -\left(\frac{d}{d y} A_{y}\right)-\frac{d}{d x} A_{x}
  \label{eq:E24}
\end{equation}

電場の発散 $\nabla\cdot\mathbf E$ を組み立てるため、$\mathbf E$ の各成分を対応する座標で偏微分する。

E11 を x で偏微分して

\begin{equation}
  \frac{\partial E_x}{\partial x} = -\left(\frac{d^2}{d x^2} \phi\right)-\frac{d^2}{d t d x} A_{x}
  \label{eq:E25}
\end{equation}

E12 を y で偏微分して

\begin{equation}
  \frac{\partial E_y}{\partial y} = -\left(\frac{d^2}{d y^2} \phi\right)-\frac{d^2}{d t d y} A_{y}
  \label{eq:E26}
\end{equation}

E13 を z で偏微分して

\begin{equation}
  \frac{\partial E_z}{\partial z} = -\left(\frac{d^2}{d z^2} \phi\right)-\frac{d^2}{d t d z} A_{z}
  \label{eq:E27}
\end{equation}

上の3つの偏微分の和として、電場の発散 $\nabla\cdot\mathbf E$ を組み立てる。

$divE = E25 + E26 + E27$

\begin{equation}
  divE = -\left(\frac{d^2}{d z^2} \phi\right)-\frac{d^2}{d y^2} \phi-\frac{d^2}{d x^2} \phi-\frac{d^2}{d t d z} A_{z}-\frac{d^2}{d t d y} A_{y}-\frac{d^2}{d t d x} A_{x}
  \label{eq:E28}
\end{equation}

組み立てた $\nabla\cdot\mathbf E$ の式を整理する。

E28 の右辺を整理して

\begin{equation}
  divE = -\left(\frac{d^2}{d z^2} \phi\right)-\frac{d^2}{d y^2} \phi-\frac{d^2}{d x^2} \phi-\frac{d^2}{d t d z} A_{z}-\frac{d^2}{d t d y} A_{y}-\frac{d^2}{d t d x} A_{x}
  \label{eq:E29}
\end{equation}

ゲージ条件 $\nabla\cdot\mathbf A=0$ の時間微分を計算する（恒等的にゼロのはずであることを確認するため）。

E20 を t で偏微分して

\begin{equation}
  \frac{\partial gauge}{\partial t} = \frac{d^2}{d t d z} A_{z}+\frac{d^2}{d t d y} A_{y}+\frac{d^2}{d t d x} A_{x}
  \label{eq:E30}
\end{equation}

この時間微分の値を 0 と置いた方程式として確定する。

E30 の右辺を 0 とする

\begin{equation}
  \frac{d^2}{d t d z} A_{z}+\frac{d^2}{d t d y} A_{y}+\frac{d^2}{d t d x} A_{x} = 0
  \label{eq:E31}
\end{equation}

ゲージ条件の時間微分＝0 を $\nabla\cdot\mathbf E$ の展開式へ代入すると、対応する3項が消えて $\nabla\cdot\mathbf E=-\nabla^2\phi$（Poisson 方程式）が得られる。

E29 に E31 を代入して

\begin{equation}
  divE = -\left(\frac{d^2}{d z^2} \phi\right)-\frac{d^2}{d y^2} \phi-\frac{d^2}{d x^2} \phi
  \label{eq:E32}
\end{equation}

\end{document}